% Copyright (c) 2026 by C. Klukas.
% Licensed under the MIT License. See LICENSE for details.

This appendix lists the hooks of \code{paglets::Paglet}, how each one
reaches the \paglet{}, and what the host promises about the order of hooks
and messages. Section~\ref{sec:lifecycle} explains the lifecycle; this
appendix only states it. Everything here is taken from
\code{cpp/sdk/include/paglets/paglet.hpp}, \code{cpp/sdk/src/paglet.cpp},
\code{cpp/common/include/paglets/abi.hpp} and the delivery loop in
\code{cpp/host/src/runtime.cpp}; the ABI specification
(\code{planning/cpp-abi-v1.md}) is the normative text for the binary
interface, summarized in Appendix~\ref{app:abi}.

\subsection{The hooks}\label{sec:events-hooks}

Each module defines one class derived from \code{paglets::Paglet} and
registers it with the macro \code{PAGLETS\_PAGLET}; the SDK keeps one static
instance of that class in linear memory. Its constructor runs inside the
first call of \code{paglets\_on\_event} or \code{paglets\_on\_message},
which for a new \paglet{} is the \code{created} event (the host's earlier
calls, \code{\_initialize} at instantiation and \code{paglets\_alloc}
before each delivery, do not touch it); a clone and an arriving \paglet{}
start from a memory image in which the object already exists, so the
constructor does not run again for them. Table~\ref{tab:events-hooks} lists
the virtual functions of \code{Paglet}. Seven of them are lifecycle events,
which the host delivers through the export \code{paglets\_on\_event} with
the event kind of section~6 of the ABI specification (``ABI section''
below names a section of \code{planning/cpp-abi-v1.md}). Three are reports, which the host delivers
as ordinary messages through \code{paglets\_on\_message}: their names start
with the reserved prefix \code{paglets.}, they carry no sender record, and
they have priority~7, the highest a message can have. The SDK recognizes
them by name and by the missing sender record, and calls the hook instead
of \code{on\_message}.
\code{on\_message} receives every other message, except a reply whose
request was given a continuation: the SDK keeps continuations in a map in
linear memory, keyed by correlation ID, and calls the continuation directly.

\begin{table}[htbp]
\centering\small
\caption{The hooks of \code{paglets::Paglet}. Events reach the \paglet{}
through \code{paglets\_on\_event} (``event~$n$'', the kind of ABI
section~6); reports are messages with priority~7 and no sender record,
recognized by name in the SDK. ``Nothing'' means the default
implementation is empty.}
\label{tab:events-hooks}
\newcommand{\rr}{\raggedright\let\\\tabularnewline}%
\setlength{\tabcolsep}{4pt}%
\begin{tabular}{@{}p{2.9cm}p{6.0cm}p{4.1cm}p{1.9cm}@{}}
\toprule
Hook & Runs & Document & Default \\
\midrule
\code{on\_created} & \rr once, as the first delivery of a new \paglet{}, before any queued message (event~1) & \rr \code{Started}: \code{args}, \code{caps}; \code{original} empty & nothing \\
\code{on\_cloned} & \rr in a clone, instead of \code{on\_created}, as its first delivery (event~5); also for a clone made on another host & \rr \code{Started}: \code{args}, \code{caps}, \code{original} (the original's ID) & nothing \\
\code{on\_activated} & \rr after the \paglet{} was restored from its memory image, right before whatever woke it: a message, a timer, a wake after \code{deactivate(ms)}, a report, or a host-side dispose or dispatch of an inactive \paglet{}; never before \code{on\_created}, \code{on\_cloned} or \code{on\_arrived} (event~2) & none & nothing \\
\code{on\_deactivating} & \rr after the handler that called \code{deactivate}, or after a host-side deactivation, before the image is taken (event~3) & none & nothing \\
\code{on\_arrived} & \rr on the new host after a move, as the first delivery there (event~4); with \code{arrival=inactive} it waits until the first message & \rr \code{abi::ArrivedEvent}: \code{from} (the previous host's key ID), \code{lost} (capabilities not re-created) & nothing \\
\code{on\_dispatching} & \rr after the handler that called \code{dispatch}, or after a host-side dispatch (the CLI's \code{dispatch}), before the image is taken (event~6) & \rr \code{abi::DispatchingEvent}: \code{destination} (the transfer ticket) & nothing \\
\code{on\_disposing} & \rr after the handler that called \code{dispose}, or after a host-side disposal, before the \paglet{} ends (event~7) & none & nothing \\
\code{on\_undelivered} & \rr an asynchronous send or reply could not be delivered (report \code{paglets.undelivered}) & \rr \code{abi::Undelivered}: \code{name}, \code{handle}, \code{status} & \rr logs a warning \\
\code{on\_move\_failed} & \rr a dispatch, or a clone for another host, did not happen; the \paglet{} continues on its host (report \code{paglets.move\_failed}) & \rr \code{abi::MoveFailed}: \code{destination}, \code{reason}, \code{clone} (empty for a dispatch) & \rr logs a warning \\
\code{on\_recovered} & \rr the host restarted and the \paglet{} came back from its last image; only after \code{set\_restart(true, true)} (report \code{paglets.recovered}, ABI~1.2) & \rr \code{abi::Recovered}: \code{image\_ms}, \code{restarted\_ms} & nothing \\
\code{on\_message} & \rr every message, request, timer, and every reply without a continuation (\code{paglets\_on\_message}) & \code{Message} & the router \\
\bottomrule
\end{tabular}
\end{table}

Appendix~\ref{app:abi} lists the handler results. Two SDK details matter
here: the SDK answers a handled request that was neither answered nor
deferred with \code{gone}, by dropping its reply capability; and the event
hooks return \code{void}, so the SDK reports \code{handled} for them and
\code{malformed} when an event document does not decode. A negative result
from \code{paglets\_on\_event} is logged by the host and does not stop the
transition.

\subsection{Ordering guarantees}\label{sec:events-ordering}

The host promises the following. Each item names the function in the
sources that implements it.

\begin{enumerate}[leftmargin=2em]
\item \textbf{One call at a time.} A \paglet{} never sees two calls at once:
  hooks and handlers run one after the other, on one thread at a time
  (\code{runtime.cpp}, \code{lane\_loop}: \code{rec.running}). Images are
  never taken while an import runs (ABI section~4).
\item \textbf{Mailbox order.} Messages are delivered by priority, highest
  first, and in arrival order within a priority. Replies and host reports
  have priority~7; a timer message has the priority given to \code{after}
  (\code{runtime.cpp}, the mailbox key \code{(-priority, seq)}; ABI
  section~7.1).
\item \textbf{First delivery.} \code{created}, \code{cloned} and
  \code{arrived} are queued above every message priority, so they are the
  \paglet{}'s first call on that host, ahead of every message queued for it
  earlier (host-side lifecycle requests share this priority and keep their
  arrival order) (\code{runtime.cpp}, \code{control\_priority}).
\item \textbf{Lifecycle operations are deferred.} \code{dispose},
  \code{deactivate}, \code{dispatch} and \code{clone} take effect when the
  current call returns. At most one of
  \code{dispose}, \code{deactivate} and \code{dispatch} can be requested
  per call; a second one fails with \code{bad\_state}. \code{clone} and
  \code{create\_child} together can be called up to
  \code{spawn\_limit\_per\_call} times (16) per call; beyond that they fail
  with \code{quota} (\code{runtime.cpp}, \code{Imports::lifecycle},
  \code{Imports::create\_child}).
\item \textbf{After the call: clones first.} When the call returns, the
  host takes one memory image for all clones requested in it and starts
  each clone with \code{cloned}; a clone for another host departs with
  that image. Only then does it run the end-of-call hook of item~6. The
  clone therefore sees the memory as the handler left it, not as
  \code{on\_dispatching} or \code{on\_deactivating} changed it
  (\code{runtime.cpp}, \code{finish\_call}).
\item \textbf{After the call: the end-of-call hook, then the image.}
  For a disposal, \code{on\_disposing} runs and the \paglet{} ends. For a
  dispatch, \code{on\_dispatching} runs, the image is taken and the move
  starts. For a deactivation, \code{on\_deactivating} runs, the image is
  taken and the instance is unloaded. What a hook changes is in the image.
  Images are taken only
  between calls, never while a call runs (\code{finish\_call}; ABI
  section~4).
\item \textbf{Checkpoints.} If none of the three operations was requested,
  the host stores an image after the call when \code{checkpoint} was called
  in it, when the \paglet{}'s checkpoint interval (default 1\,s) has
  elapsed, after the first call of a \paglet{} that has no stored image,
  and after \code{set\_restart}. Nothing is stored without a state
  directory or with \code{restart} false (\code{finish\_call},
  \code{persist}).
\item \textbf{Activation.} A \paglet{} without a loaded instance is restored
  from its image before the next delivery, a host-side dispose or dispatch
  included, and \code{on\_activated} runs right before that delivery. It
  does not run before \code{created}, \code{cloned} or \code{arrived}: a new
  instance, a clone and an arriving \paglet{} get their start event only
  (\code{runtime.cpp}, \code{process}, \code{activate}).
\item \textbf{A pending move holds messages.} From the image until the
  mesh commits the move or reports it failed, the \paglet{} is unloaded and
  messages for it are held on the old host. After a successful move they
  are forwarded and delivered after \code{on\_arrived}, in mailbox order
  with whatever reached the new host directly. After a failed move the
  \paglet{} is restored on the old host: \code{on\_activated}, then the held
  messages and the report in mailbox order, where \code{paglets.move\_failed}
  is queued at priority~7 when the outcome is known, so it precedes every
  held message of lower priority and follows held replies and held messages
  sent at priority~7 (\code{Runtime::finish\_departure}).
\item \textbf{Timers and wakes.} A timer or a wake after
  \code{deactivate(ms)} activates a deactivated \paglet{}; a wake alone
  delivers no message, only \code{on\_activated}. Timers survive
  deactivation, moves and host restarts (\code{Impl::fire}, \code{recover}).
  A wake does not: it is not stored with the \paglet{} and does not travel
  with it, so after a host restart (or a host-side dispatch of the sleeping
  \paglet{}) the \paglet{} waits for its next message or timer.
\item \textbf{Recovery.} After a host restart a \paglet{} with
  \code{restart} true is re-created from its last image with its
  capabilities, timers and the correlation IDs of its pending requests.
  Those requests are answered \code{timeout}, and the answers are queued
  before \code{paglets.recovered}; the first calls are then
  \code{on\_activated}, the timeout replies, \code{on\_recovered} (when
  requested), and the rest of the mailbox (\code{recover}).
\item \textbf{Budget and failure.} Every call, hooks included, has the
  handler time budget (default 5\,s). A trap or an exceeded budget in a
  message handler or in \code{on\_created}, \code{on\_cloned},
  \code{on\_arrived}, \code{on\_activated} or \code{on\_deactivating} ends the
  \paglet{} as failed; pending requests to it are answered \code{failed}. A
  trap in \code{on\_disposing} is logged and the \paglet{} ends as disposed.
  A trap in \code{on\_dispatching} is reported to the \paglet{} as a failed
  move (\code{reason} ``no image: \ldots'') and it stays
  (\code{guest\_call}, \code{process}, \code{finish\_call}).
\end{enumerate}

\subsection{The sequences}\label{sec:events-sequences}

Figure~\ref{fig:events-sequences} draws items 4 to 6 and 8 to 10 for the
four operations. Each row starts with the handler call that requested the
operation and ends with the next call the \paglet{} receives. White nodes
are calls into the \paglet{}; shaded nodes are steps of the host; the grey
labels without a frame mark what happens at that point.

\begin{figure}[htbp]
\centering
\begin{tikzpicture}[
  node distance=2mm and 1.3mm,
  plain/.append style={font=\scriptsize,inner xsep=3.5pt,inner ysep=3.5pt},
  box/.append style={font=\scriptsize,inner xsep=3.5pt,inner ysep=3.5pt},
  rowlabel/.style={font=\scriptsize\bfseries\color{pagletsBlue},anchor=east,text width=12.5mm,align=right},
  every edge/.style={arrow}
]
% -- dispatch -------------------------------------------------------------
\node[rowlabel] (l1) at (0,0) {dispatch};
\node[plain,right=1.5mm of l1] (d1) {handler calls\\\code{dispatch(t)}};
\node[plain,right=of d1] (d2) {\code{on\_dispatching}\\\code{\{destination\}}};
\node[box,right=of d2] (d3) {image};
\node[box,right=of d3] (d4) {transfer\\to host B};
\node[plain,right=of d4] (d5) {\code{on\_arrived}\\on host B};
\node[plain,right=of d5] (d6) {next delivery\\on host B};
\draw (d1) edge (d2) (d2) edge (d3) (d3) edge (d4) (d4) edge (d5) (d5) edge (d6);
\node[event,below=0.8mm of d1] {returns};
\node[event,below=0.8mm of d3] {messages\\held on A};
% failed move
\node[plain,below=8.5mm of d4.south,anchor=north west,xshift=-7mm] (f1) {\code{on\_activated}\\on host A};
\node[plain,right=of f1] (f2) {\code{on\_move\_failed}\\\code{\{destination, reason\}}};
\node[plain,right=of f2] (f3) {held messages\\(replies first)};
\draw[dashedarrow] (d4.south) -- (d4.south |- f1.north);
\node[event,anchor=west] at ($(d4.south)+(0.8mm,-4.5mm)$) {fails};
\draw (f1) edge (f2) (f2) edge (f3);
% -- deactivate -----------------------------------------------------------
\node[rowlabel,below=21mm of l1] (l2) {deactivate};
\node[plain,right=1.5mm of l2] (a1) {handler calls\\\code{deactivate(ms)}};
\node[plain,right=of a1] (a2) {\code{on\_deactivating}};
\node[box,right=of a2] (a3) {image,\\unload};
\node[box,right=of a3] (a4) {message,\\timer or wake};
\node[box,right=of a4] (a5) {restore\\from image};
\node[plain,right=of a5] (a6) {\code{on\_activated}};
\node[plain,right=of a6] (a7) {the message};
\draw (a1) edge (a2) (a2) edge (a3) (a3) edge[dashedarrow] (a4) (a4) edge (a5) (a5) edge (a6) (a6) edge (a7);
% -- dispose --------------------------------------------------------------
\node[rowlabel,below=9mm of l2] (l3) {dispose};
\node[plain,right=1.5mm of l3] (p1) {handler calls\\\code{dispose()}};
\node[plain,right=of p1] (p2) {\code{on\_disposing}};
\node[box,right=of p2] (p3) {ended: queued requests and held\\reply capabilities answered \code{gone}};
\draw (p1) edge (p2) (p2) edge (p3);
% -- clone ----------------------------------------------------------------
\node[rowlabel,below=12mm of l3] (l4) {clone};
\node[plain,right=1.5mm of l4] (c1) {handler calls\\\code{clone\_raw(args)}};
\node[box,right=of c1] (c2) {image};
\node[plain,right=of c2] (c3) {original: end-of-call\\hook, next delivery};
\node[plain,below=2mm of c3.south west,anchor=north west] (c4) {clone, here or on host B: \code{on\_cloned}\\\code{\{original, args, caps\}}, next delivery};
\draw (c1) edge (c2) (c2) edge (c3);
\draw[arrow] (c2.east) -- ++(1mm,0) |- (c4.west);
\end{tikzpicture}
\caption{The order of calls and host steps after a handler requests a
lifecycle operation (items 4 to 6 and 8 to 10 of
Section~\ref{sec:events-ordering}). The dashed path of the first row is a
move that the mesh reports as failed; the dashed arrow of the second row is
the wait, of any length, for the event that activates the \paglet{} again.}
\label{fig:events-sequences}
\end{figure}

Two consequences are easy to miss. First, as item~8 says, a move delivers
\code{on\_arrived} and no \code{on\_activated} on the new host; in the
course of a move, the only \code{on\_activated} a \paglet{} can see is the
one on its old host after a failed move, or the one before
\code{on\_dispatching} when a host-side dispatch wakes an inactive
\paglet{}. Second, the continuation of a request issued before
\code{dispatch} lives in linear memory and travels with the \paglet{}; the
reply follows it (Section~\ref{sec:communication}), is held while the move
is pending, and is delivered after \code{on\_arrived} like any other
message.
